% ==============================================================================
% 1. 基础设置与字体
% ==============================================================================
\usepackage{xeCJK}
\xeCJKsetup{
    PunctStyle=kaiming,                  % 开明式：句末标点始终压缩为半角，腾出更多空间
    CheckSingle=true                     % 避免段末单字独占一行
}
\tolerance=1000                          % 提高断行容忍度，允许稍松的行以避免标点溢出到下一行
\emergencystretch=3em                    % 加大紧急伸缩量
% --- CJK 字体：自适应选择 ---
% Windows 使用 SimSun / SimHei / KaiTi；macOS 等系统缺少这些字体时，
% 回退到对应的系统中文字体（Songti SC / STHeiti / Kaiti SC 等）。
\ExplSyntaxOn
\cs_new_protected:Npn \exam_first_font:Nn #1#2 {
  \tl_gclear:N #1
  \clist_map_inline:nn {#2} {
    \IfFontExistsTF {##1} { \tl_gset:Nn #1 {##1} \clist_map_break: } { }
  }
}
\NewDocumentCommand{\examSetCJKMain}{}{
  \exam_first_font:Nn \l_exam_cjk_main   {SimSun, Songti~SC, STSong}
  \exam_first_font:Nn \l_exam_cjk_bold   {SimHei, STHeiti, Heiti~SC}
  \exam_first_font:Nn \l_exam_cjk_italic {KaiTi, Kaiti~SC, STKaiti}
  \tl_if_empty:NT \l_exam_cjk_italic {
    \tl_gset:Nx \l_exam_cjk_italic { \l_exam_cjk_main }
  }
  \setCJKmainfont{\l_exam_cjk_main}[
    BoldFont = \l_exam_cjk_bold,
    ItalicFont = \l_exam_cjk_italic
  ]
}
\ExplSyntaxOff
\examSetCJKMain
\setmainfont{Times New Roman}

% 全局段落设置：无缩进，无段间距，正文行距
\setlength{\parindent}{0pt}
\setlength{\parskip}{0pt}
\renewcommand{\baselinestretch}{1.5}

% Keep a finite amount of stretch on prose paragraph final lines.  The
% default infinite \parfillskip makes a one-character final line cost-free,
% so TeX may strand a single Han character or punctuation mark.  A small
% finite reserve (half the local line width) makes that break expensive while
% preserving ordinary justification.  Structured layouts (TOC, choices, figures and tables)
% restore their own paragraph settings or never enter this hook.
\newcommand{\examproseparagraph}{%
    \parfillskip=0pt plus 0.5\linewidth
}

% 单字不成行、单行不成页
\makeatletter
\newif\ifexamcolorstackinitialized
\ifdefined\main@pdfcolorstack
    \examcolorstackinitializedtrue
\else
    \examcolorstackinitializedfalse
\fi
\newcommand{\examwidoworphanrules}{%
    \clubpenalty=10000%
    \widowpenalty=10000%
    \displaywidowpenalty=10000%
    \brokenpenalty=10000%
    \@clubpenalty=10000%
}
\makeatother
\examwidoworphanrules
\AtBeginDocument{\examwidoworphanrules}
\raggedbottom

% ==============================================================================
% 2. 章节标题样式 (试卷结构)
% ==============================================================================
\ctexset {
    chapter = {
        pagestyle = headings,
        break = \clearpage,      % 关键：每一章(每一份试卷)强制换页
        numbering = false,       % 关键：不显示“第一章”字样，仅显示标题
        format = \bfseries\zihao{3}\centering,  % 三号字
        beforeskip = 0em,
        afterskip = 0em,
    },
    section = {
        format = \bfseries\zihao{-4}, % 小四号字并加粗
        number = \chinese{section}、,  % 格式：一、
        indent = 0pt,
        aftername = \hspace{0pt},
        beforeskip = 0.5em,
        afterskip = 0.5em,
    },
}

\usepackage{etoolbox}
% tall-page 布尔开关：\tallpagetrue（默认）/ \tallpagefalse
\newif\iftallpage
\tallpagetrue
\makeatletter
% Reset the TeX/xcolor current-color bookkeeping without pushing a new
% graphics color-stack entry.  Chapter hooks run at top level, so the usual
% \normalcolor (which calls \set@color and schedules an aftergroup pop) would
% leak one stack entry per chapter.
\newcommand{\examresetnormalcolor}{%
    \let\current@color\default@color
    \ifdefined\XC@default@color
        \let\XC@current@color\XC@default@color
    \fi
    \ifexamcolorstackinitialized\else
        \special{pdfcolorstackinit 1 page direct (0 g 0 G)}%
        \global\examcolorstackinitializedtrue
    \fi
    \ifdefined\@pdfcolorstack
        \special{pdfcolorstack \@pdfcolorstack set (0 g 0 G)}%
    \else
        \special{pdfcolorstack 1 set (0 g 0 G)}%
        % Legacy xdv color handlers still understand the direct `color`
        % special; discard one leaked non-default color first, then change
        % the graphics state in place without another stack push.
        \ifx\current@color\default@color\else
            \special{color pop}%
        \fi
        \special{color gray 0}%
    \fi
}
\makeatother
\pretocmd{\chapter}{%
    \FloatBarrier
    \examresetnormalcolor
    \problembodyend
    \setcounter{probnum}{0}%
    \setcounter{section}{0}%
    \iftallpage
        \pdfpageheight=550cm\relax
        \setlength{\textheight}{546cm}%
        \vsize=\textheight\relax
    \fi
}{}{}
\pretocmd{\section}{\problembodyend}{}{}
\makeatletter
\patchcmd{\tableofcontents}
  {\@starttoc{toc}}
  {%
    \iftallpage
      \@starttoc{toc}%
    \else
      \begingroup
        \normalfont
        \zihao{-4}\selectfont
        \setlength{\columnsep}{1.8em}%
        \setlength{\columnseprule}{0.8pt}%
        \begin{multicols}{2}\@starttoc{toc}\end{multicols}%
      \endgroup
    \fi
  }
  {}{}
\renewcommand*\l@part[2]{%
  \addpenalty{-\@highpenalty}%
  \vskip 0.8em \@plus\p@
  \begingroup
    \parindent \z@
    \rightskip \z@
    \parfillskip \z@
    \centering\normalfont #1\par
  \endgroup
}
\renewcommand*\l@chapter[2]{%
  \ifnum \c@tocdepth >\m@ne
    \addpenalty{\@secpenalty}%
    \addvspace{0pt}%
    \begingroup
      \parindent \z@
      \rightskip \@pnumwidth
      \parfillskip -\@pnumwidth
      \leavevmode
      #1\nobreak\leaders\hbox{$\m@th \mkern 2mu.\mkern 2mu$}\hfill\nobreak
      \hb@xt@\@pnumwidth{\hss #2}\par
    \endgroup
  \fi
}
\makeatother

% 正文小四号字
\AtBeginDocument{\zihao{-4}}

\newcommand{\examspecialsection}[1]{%
    \problembodyend
    \par
    \medskip
    \begin{center}
    \bfseries\zihao{-4}#1
    \end{center}
    \setcounter{section}{0}%
    \medskip
}
\newcommand{\additionalproblemsection}{\examspecialsection{附加题}}
\newcommand{\referenceproblemsection}{\examspecialsection{参考题}}
\newcommand{\optionalproblemsection}{\examspecialsection{选作题}}

% ==============================================================================
% 3. 宏包加载 (精简版)
% ==============================================================================
\usepackage{geometry}
\geometry{a4paper,left=2cm,right=2cm,top=2cm,bottom=2cm}
\newif\ifexampageflexused
\AddToHook{shipout/after}{\global\exampageflexusedfalse}
\newcommand{\examavoidshortpage}{%
    \ifexampageflexused\else
        \global\exampageflexusedtrue
        \relax
    \fi
}

% 单栏排版；目录页局部使用 multicol
\usepackage{multicol}

% 数学相关
\usepackage{amsmath, amssymb, amsthm, bm}
\usepackage{unicode-math}
% 数学字体：默认由构建脚本 (tools/texlive-font-paths.sh) 通过绝对路径加载，
% 跨平台且对 macOS 友好（macOS 上 XeTeX 走 Core Text，按字体名找不到 TeX Live 的数学字体）。
% 若直接 xelatex（未走 make），则回退到按字体名加载（Windows / 标准 TeX Live 可用）。
\InputIfFileExists{tmp/math-fonts}{}{%
    \setmathfont{TeX Gyre Termes Math}%
    % 希腊字母使用 STIX 字体 (MathType 风格，正体)
    \setmathfont{STIX}[
        range = {up/{Greek,greek}, it/{Greek,greek}, bfup/{Greek,greek}, bfit/{Greek,greek}}
    ]%
    % 大运算符和空集使用 Latin Modern Math (LaTeX 默认风格)
    \setmathfont{Latin Modern Math}[
        range = {\sum, \prod, \coprod, \int, \iint, \iiint, \iiiint, \oint, \bigcup, \bigcap, \bigsqcup, \bigvee, \bigwedge, \bigoplus, \bigotimes, \bigodot}
    ]%
    \setmathfont{Asana Math}[
        range = {\varnothing}
    ]%
}
\let\emptyset\varnothing

\allowdisplaybreaks

% cases 行距与正文节奏保持一致；amsmath 的 cases 内部会重设 \arraystretch，
% 所以直接重定义内部入口，避免外层 hook 被覆盖。
\makeatletter
\def\env@cases{%
  \let\@ifnextchar\new@ifnextchar
  \left\lbrace
  \def\arraystretch{0.9}%
  \array{@{}l@{\enspace}l@{}}%
}
% An inline cases block is one math atom.  Discard source-line whitespace
% immediately after the environment so a following Chinese clause is not
% separated by an accidental interword glue node.
\def\endcases{\endarray\right.\ignorespacesafterend}
\makeatother

% 图形与表格
\usepackage{graphicx}
\usepackage{adjustbox}
\usepackage{xparse}

\usepackage{tikz}
\usetikzlibrary{arrows.meta, calc, intersections, positioning, shapes.geometric, shapes.misc}
\usepackage{pgfplots}
\pgfplotsset{compat=1.18}
\pgfplotsset{
    ybar interval/.append style={
        x tick label as interval=false,
    },
}
\input{statistical-charts.tex}

\usepackage{wrapfig}
% 调整 wrapfigure 与正文的间距；普通题内 figure 也采用同一稳定半行距。
\setlength{\intextsep}{0.45\baselineskip}  % 纵向：图片上下的间距
\setlength{\textfloatsep}{0.45\baselineskip}
\setlength{\floatsep}{0.45\baselineskip}
% \columnsep（横向：图片与侧边文字的间距）已在上方 multicol 区域统一设为 1.5em
\usepackage{array, tabularray, multirow}
% 正文依靠 \lineskip 在含显示样式分式的相邻行之间自动避让。array/tabular
% 内部会把 \lineskip 清零，故进入表格前保存正文间距的一半，由分式上下各承担
% 一半；相邻两个高行合计仍是一个正文 \lineskip，普通文本行不被无谓拉高。
\newif\ifexamtablefractioncontext
\newdimen\examtablefractionclearance
\newcommand{\examcapturetablefractionclearance}{%
    \examtablefractionclearance=\lineskip
    \divide\examtablefractionclearance by 2
}
\usepackage{caption}
\usepackage{placeins}
% Permit large block figures to use the standard top-float queue even when
% their measured box is taller than the conservative defaults.
\renewcommand{\topfraction}{0.95}
\renewcommand{\textfraction}{0.05}
\renewcommand{\floatpagefraction}{0.70}
\setcounter{topnumber}{2}
% Keep float-only pages compact and top-aligned.  The default \@fptop and
% \@fpbot both contain stretch, which vertically centers a lone figure;
% reserve the stretch only below the float instead.
\makeatletter
\setlength{\@fptop}{0pt}
\setlength{\@fpsep}{0.5\baselineskip}
\setlength{\@fpbot}{0pt plus 1fil}
\makeatother
\AtEndDocument{\FloatBarrier}

% --- 接管 \includegraphics：找不到文件时显示占位框而非报错 ---
\let\originalincludegraphics\includegraphics
\renewcommand{\includegraphics}[2][]{%
    \problemconsumeinlineblockprefix
    \IfFileExists{#2}{%
        \originalincludegraphics[#1]{#2}%
    }{%
        \setlength{\fboxsep}{4pt}%
        \fbox{%
            \begin{minipage}[c][4em][c]{0.8\linewidth}%
                \centering\color{gray}\small 图片未找到：{\ttfamily\detokenize{#2}}%
            \end{minipage}%
        }%
    }%
}

\ExplSyntaxOn
\bool_new:N \l__bitmap_solution_caption_bool
\tl_new:N \l__bitmap_solution_caption_tl

\cs_new_protected:Npn \__bitmap_prepare_caption:n #1
{
    \bool_set_false:N \l__bitmap_solution_caption_bool
    \tl_clear:N \l__bitmap_solution_caption_tl
    \regex_extract_once:nnNT { q([0-9]+)_solution_fig([0-9]+)\.(?:tex|png|jpg|jpeg)$ } {#1} \l_tmpa_seq
    {
        \bool_set_true:N \l__bitmap_solution_caption_bool
        \tl_set:Nx \l__bitmap_solution_caption_tl
        {第\seq_item:Nn \l_tmpa_seq {2}题解析图\seq_item:Nn \l_tmpa_seq {3}}
    }
}
\NewDocumentCommand{\bitmappreparecaption}{m}
{
    \__bitmap_prepare_caption:n {#1}
}
\NewDocumentCommand{\bitmapinsertcaption}{}
{
    \bool_if:NTF \l__bitmap_solution_caption_bool
        {\captionof*{figure}{\tl_use:N \l__bitmap_solution_caption_tl}}
        {\captionof{figure}{}}
}
\ExplSyntaxOff

\makeatletter
\newcommand{\FigureLayoutDeclare}[3]{%
    \expandafter\gdef
        \csname bitmaplayout@\detokenize{#1}\endcsname
        {trim={#3},clip,width=#2,max width=\linewidth}%
}
% Trim-only declarations preserve the caller's width/max width controls.  A
% complete generated layout remains authoritative when both declarations
% exist for the same image.
\newcommand{\FigureTrimDeclare}[2]{%
    \expandafter\gdef
        \csname bitmaptrim@\detokenize{#1}\endcsname
        {#2}%
}
\InputIfFileExists{tools/typeset/generated_image_layout.tex}{}{}
\InputIfFileExists{tools/crop/manual_image_layout.tex}{}{}

\newcommand{\bitmapinclude}[2][]{%
    \begingroup
    \ifcsname bitmaplayout@\detokenize{#2}\endcsname
        % 电子稿位图采用生成的白边裁切与标签定标宽度；
        % 调用处旧的 width/height/scale 不再参与，避免同类题图忽大忽小。
        \edef\bitmaplayoutoptions{%
            \csname bitmaplayout@\detokenize{#2}\endcsname
        }%
    \else\ifcsname bitmaptrim@\detokenize{#2}\endcsname
        % Trim-only assets keep the invocation's width/max width while
        % applying the declared crop.
        \edef\bitmaplayoutoptions{%
            trim={\csname bitmaptrim@\detokenize{#2}\endcsname},clip,max width=\linewidth,#1%
        }%
    \else
        % 照片、扫描稿、选项图及低置信度图片保留原调用参数。
        \def\bitmaplayoutoptions{max width=\linewidth,#1}%
    \fi\fi
    \expandafter\adjustbox\expandafter{\bitmaplayoutoptions}{\includegraphics{#2}}%
    \endgroup
}
\makeatother

% TikZ/TeX figure bodies may be read while a problem's continuation
% \everypar is active.  Keep that paragraph layout out of the boxed figure;
% the surrounding problem paragraph remains unchanged by this local group.
\newcommand{\texboxparagraphreset}{%
    \everypar{}%
    \leftskip=0pt\relax
    \rightskip=0pt\relax
    \parindent=0pt\relax
    \hangindent=0pt\relax
    \hangafter=1\relax
    \parshape=0\relax
}

% TikZ's /tikz/font key is evaluated after the node text hbox has started in
% this project, which can leave an invisible leading width when it contains a
% font command (for example, \normalsize).  Redirect it locally to /tikz/node
% font so the font is selected before TikZ measures the text hbox.  The key
% definition lives only inside each figure wrapper group.
\newcommand{\texboxtikzfontredirect}{%
    \pgfkeys{/tikz/font/.code={\pgfkeysalso{/tikz/node font={##1}}}}%
}

% TeX figures are rendered at their authored size, then the finished box is
% uniformly scaled.  Scaling the box (rather than TikZ coordinates) keeps
% lines, anchors, nodes, labels and PGFPlots ticks in exactly the same relative
% geometry.  Source-level \small/\scriptsize/\tiny declarations therefore
% remain meaningful and are not allowed to create a second, shape-only scale.
\makeatletter
\newsavebox{\texfiguremeasurebox}
\newlength{\texfigurenaturalwidth}
\newlength{\texfiguretargetwidth}
\newlength{\texfiguremaxwidth}
\newlength{\texfigurefinalwidth}
\newlength{\texfiguremaxnodefontdim}
\newlength{\texfigurenodefontdim}
\newlength{\texfigureautotargetwidth}
\newlength{\texfigurebboxmargin}
\newlength{\texfiguresavedmaxnodefontdim}
\setlength{\texfigurebboxmargin}{2.5pt}
\def\texfigure@width{}
\def\texfigure@maxwidth{}
\def\texfigure@scale{1}
\def\texfigure@haswidth{0}
\def\texfigure@hasmaxwidth{0}
\def\texfigure@factor{1}
\def\texfigure@fontfactor{1}
\def\texfigure@bodyfontsize{10}

\def\texfigure@parseone#1=#2\@nil{%
    \edef\texfigure@key{\zap@space#1 \@empty}%
    \def\texfigure@value{#2}%
    \ifdefstring{\texfigure@key}{width}{%
        \def\texfigure@width{#2}%
        \def\texfigure@haswidth{1}%
    }{}%
    % Spaces are removed from keys above, so `max width` is compared as
    % `maxwidth` here.
    \ifdefstring{\texfigure@key}{maxwidth}{%
        \def\texfigure@maxwidth{#2}%
        \def\texfigure@hasmaxwidth{1}%
    }{}%
    \ifdefstring{\texfigure@key}{scale}{\def\texfigure@scale{#2}}{}%
}
\def\texfigure@parseopts#1{%
    \def\texfigure@width{}%
    \def\texfigure@maxwidth{}%
    \def\texfigure@scale{1}%
    \def\texfigure@haswidth{0}%
    \def\texfigure@hasmaxwidth{0}%
    \@for\texfigure@item:=#1\do{%
        \expandafter\texfigure@parseone\texfigure@item\@nil
    }%
}

\def\texfigure@input#1{%
    \texboxparagraphreset
    \texboxtikzfontredirect
    % Legacy sources often contain a generous `use as bounding box`
    % rectangle.  Keep that path visible but make it overlay-only, so TikZ's
    % automatic box follows actual strokes and labels.  The enclosing measure
    % adds a small fixed margin.
    \pgfkeys{/tikz/use as bounding box/.style={overlay}}%
    \input{#1}%
}

% TikZ and PGFPlots execute this hook after the node's font keys have been
% applied.  The assignment is global only for the duration of one measurement:
% texinclude snapshots and restores the register below, so nested figure
% groups or surrounding problem text cannot observe the measured value.
\def\texfigure@recordnodefont{%
    \ifx\f@size\@undefined\else
        \texfigurenodefontdim=\f@size pt\relax
        \ifdim\texfigurenodefontdim>\texfiguremaxnodefontdim
            \global\texfiguremaxnodefontdim=\texfigurenodefontdim
        \fi
    \fi
}

\def\texfigure@installrecordhooks{%
    % Append directly to TikZ's internal begin-node accumulator.  Unlike an
    % `every node/.append style` entry, this survives a source picture that
    % later replaces `every node/.style=...`; source execute-at-begin-node code
    % is itself append-only and therefore cannot erase this recorder.
    \expandafter\def\expandafter\tikz@atbegin@node\expandafter{\tikz@atbegin@node\texfigure@recordnodefont}%
    % A few legacy nodes select their font from an execute-at-begin hook.  A
    % matching end hook samples the resulting font after such source code has
    % run, while remaining inside TikZ's node text group.
    \expandafter\def\expandafter\tikz@atend@node\expandafter{\tikz@atend@node\texfigure@recordnodefont}%
    \pgfkeys{/tikz/every node/.append style={execute at begin node={\texfigure@recordnodefont}}}%
    % PGFPlots tick, title and axis-description nodes normally inherit
    % /tikz/every node; these explicit hooks cover styles that replace it.
    \pgfkeys{/pgfplots/every tick label/.append style={execute at begin node={\texfigure@recordnodefont}}}%
    \pgfkeys{/pgfplots/every axis label/.append style={execute at begin node={\texfigure@recordnodefont}}}%
    \pgfkeys{/pgfplots/every axis title/.append style={execute at begin node={\texfigure@recordnodefont}}}%
}

\def\texfigure@measure#1{%
    \setbox\texfiguremeasurebox=\hbox{%
        \kern\texfigurebboxmargin
        \begingroup
            \texfigure@input{#1}%
        \endgroup
        \kern\texfigurebboxmargin
    }%
    % Legacy flowcharts frequently declare `baseline` at coordinate y=0
    % while all visible nodes sit above it.  Rebase the finished picture on
    % its visible bottom so the box carries no artificial depth below the
    % artwork (choice figures intentionally bypass texinclude and keep their
    % dedicated R5-4a baseline behavior).
    \ifdim\dp\texfiguremeasurebox>0pt
        \setbox\texfiguremeasurebox=\hbox{%
            \raisebox{\dp\texfiguremeasurebox}{\box\texfiguremeasurebox}%
        }%
    \fi
    \texfigurenaturalwidth=\wd\texfiguremeasurebox
}

% Determine one factor for the complete box.  Automatic enlargement raises
% the largest observed node label to the current body size.  Historical
% width/max width/scale values below that target are deliberately ignored;
% explicit enlargement remains honoured, subject only to the final line-width
% cap.  max width is an upper bound only when it is not smaller than the
% automatic font target.
\def\texfigure@setfactor{%
    \def\texfigure@fontfactor{1}%
    \ifdim\texfigurenaturalwidth>0pt
        \ifdim\texfiguremaxnodefontdim>0pt
            \pgfmathparse{max(1,\texfigure@bodyfontsize / \strip@pt\texfiguremaxnodefontdim)}%
            \edef\texfigure@fontfactor{\pgfmathresult}%
        \fi
        \edef\texfigure@factor{\texfigure@fontfactor}%
        % scale<1 is a historical shrink request; scale>1 is an explicit
        % enlargement and participates in the same uniform factor.
        \pgfmathparse{max(\texfigure@factor,\texfigure@scale)}%
        \edef\texfigure@factor{\pgfmathresult}%
        \ifnum\texfigure@haswidth=1
            \texfiguretargetwidth=\texfigure@width\relax
            \ifdim\texfiguretargetwidth>\texfigurenaturalwidth
                \edef\texfigure@ratio{\strip@pt\texfiguretargetwidth/\strip@pt\texfigurenaturalwidth}%
                \pgfmathparse{max(\texfigure@factor,\texfigure@ratio)}%
                \edef\texfigure@factor{\pgfmathresult}%
            \fi
        \fi
        % Compute the automatic target in physical units.  A max width below
        % this target is ignored, preserving readable labels from old calls.
        \pgfmathparse{\texfigure@fontfactor * \strip@pt\texfigurenaturalwidth}%
        \texfigureautotargetwidth=\pgfmathresult pt\relax
        \ifnum\texfigure@hasmaxwidth=1
            \texfiguremaxwidth=\texfigure@maxwidth\relax
            \ifdim\texfiguremaxwidth>\texfigureautotargetwidth
                \pgfmathparse{\texfigure@factor * \strip@pt\texfigurenaturalwidth}%
                \texfiguretargetwidth=\pgfmathresult pt\relax
                \ifdim\texfiguretargetwidth>\texfiguremaxwidth
                    \edef\texfigure@ratio{\strip@pt\texfiguremaxwidth/\strip@pt\texfigurenaturalwidth}%
                    \pgfmathparse{min(\texfigure@factor,\texfigure@ratio)}%
                    \edef\texfigure@factor{\pgfmathresult}%
                \fi
            \fi
        \fi
    \else
        \def\texfigure@factor{1}%
    \fi
    % Final hard cap: only the actual finished box width may force a shrink.
    \ifdim\texfigurenaturalwidth>0pt
        \pgfmathparse{\texfigure@factor * \strip@pt\texfigurenaturalwidth}%
        \texfigurefinalwidth=\pgfmathresult pt\relax
        \ifdim\texfigurefinalwidth>\linewidth
            % Leave a tiny numerical safety margin: pgfmath's decimal
            % rounding can otherwise make \scalebox one or two milli-points
            % wider than \linewidth in the emitted box.
            \texfiguretargetwidth=\dimexpr\linewidth-0.1pt\relax
            \edef\texfigure@ratio{\strip@pt\texfiguretargetwidth/\strip@pt\texfigurenaturalwidth}%
            \pgfmathparse{min(\texfigure@factor,\texfigure@ratio)}%
            \edef\texfigure@factor{\pgfmathresult}%
            \texfigurefinalwidth=\texfiguretargetwidth
        \fi
    \else
        \texfigurefinalwidth=0pt
    \fi
}

\newcommand{\texinclude}[2][]{%
    \begingroup
        \texfigure@parseopts{#1}%
        \edef\texfigure@bodyfontsize{\f@size}%
        \texfiguresavedmaxnodefontdim=\texfiguremaxnodefontdim
        \global\texfiguremaxnodefontdim=0pt\relax
        \texfigure@installrecordhooks
        \texfigure@measure{#2}%
        \texfigure@setfactor
        \typeout{TEXFIGURE body-font=\texfigure@bodyfontsize pt recorded-max-node-font=\the\texfiguremaxnodefontdim uniform-factor=\texfigure@factor natural-width=\the\texfigurenaturalwidth final-width=\the\texfigurefinalwidth}%
        \scalebox{\texfigure@factor}{\box\texfiguremeasurebox}%
        \global\texfiguremaxnodefontdim=\texfiguresavedmaxnodefontdim
    \endgroup
}
\makeatother

% Block figure float support.  The outer float key is consumed here; every
% other adjustbox key is forwarded unchanged to bitmapinclude/texinclude.
\makeatletter
\newsavebox{\examfigurebox}
\newlength{\examfigureheight}
% Remember the source page of each exam figure float.  LaTeX may otherwise
% put a deferred condition figure at the top of the page where its question
% starts, which reverses the question/figure order.  The source page is keyed
% by the float box register and therefore survives the deferlist.
\newif\ifexamfigure@recordsource
\examfigure@recordsourcefalse
\newcount\examfigure@sourcepage
\examfigure@sourcepage=0
\newdimen\examfigure@pendingheight
\examfigure@pendingheight=0pt
% Height of mapped exam figures already inserted at the current page top.
% Unlike pendingheight, this must survive queue consumption until the page is
% shipped: tcolorbox may still compute its page height after the top float has
% left the defer list but before it is physically shipped.
\newdimen\examfigure@currenttopheight
\examfigure@currenttopheight=0pt
\newcount\examfigure@currenttoppage
\examfigure@currenttoppage=-1
\newif\ifexamfigure@tcbreserve
\examfigure@tcbreservefalse
\newif\ifexamfigure@tcbready
\examfigure@tcbreadyfalse
\def\examfigure@opts{}
\def\examfigure@float{auto}
\def\examfigure@modeinline{inline}
\def\examfigure@modetop{top}
\ExplSyntaxOn
\tl_new:N \l__exam_figure_raw_tl
\tl_new:N \l__exam_figure_float_tl
\cs_new:Npn \examfigureparsedopts {}
\cs_new:Npn \examfigureparsedfloat {auto}
\cs_new_protected:Npn \examfigureparseopts #1
  {
    \tl_set:Nn \l__exam_figure_raw_tl {#1}
    \tl_set:Nn \l__exam_figure_float_tl {auto}
    \tl_if_in:VnT \l__exam_figure_raw_tl {float}
      {
        \regex_extract_once:nnNT
          { (?:\A|,)\s*float\s*=\s*(auto|inline|top)\s*(?:,|\Z) }
          { \tl_to_str:V \l__exam_figure_raw_tl } \l_tmpa_seq
          { \tl_set:Nx \l__exam_figure_float_tl { \seq_item:Nn \l_tmpa_seq {1} } }
        \regex_replace_once:nnN
          { \A\s*float\s*=\s*(?:auto|inline|top)\s*,\s* }
          { }
          \l__exam_figure_raw_tl
        \regex_replace_once:nnN
          { ,\s*float\s*=\s*(?:auto|inline|top)\s*,\s* }
          { , }
          \l__exam_figure_raw_tl
        \regex_replace_once:nnN
          { \A\s*float\s*=\s*(?:auto|inline|top)\s*\Z }
          { }
          \l__exam_figure_raw_tl
        \regex_replace_once:nnN
          { ,\s*float\s*=\s*(?:auto|inline|top)\s*\Z }
          { }
          \l__exam_figure_raw_tl
      }
    % Publish the parsed values as live macros.  A \cs_new_eq:NN alias here
    % would snapshot the initially empty token lists and silently discard
    % every bitmapfigure/texfigure option.
    \cs_set:Npx \examfigureparsedopts
      { \exp_not:V \l__exam_figure_raw_tl }
    \cs_set:Npx \examfigureparsedfloat
      { \exp_not:V \l__exam_figure_float_tl }
  }
\ExplSyntaxOff
\def\examfigure@parseopts#1{%
    \examfigureparseopts{#1}%
    % Copy one expansion level so xkeyval receives a literal comma-separated
    % key list while commands inside values (for example \linewidth) survive.
    \expandafter\def\expandafter\examfigure@opts
        \expandafter{\examfigureparsedopts}%
    \edef\examfigure@float{\examfigureparsedfloat}%
}
\def\examfigure@buildbox#1{%
    \setbox\examfigurebox=\vbox{%
        \hsize=\linewidth
        \begin{minipage}{\linewidth}
            \centering
            \begingroup
                \def\theprobnum{\examfigure@frozenprobnum}%
                #1%
                \bitmapinsertcaption
            \endgroup
        \end{minipage}%
    }%
    \examfigureheight=\dimexpr\ht\examfigurebox+\dp\examfigurebox\relax
}
% Inline figures are ordinary block content in a problem.  Keep a stable,
% compact half-line separation from surrounding prose; the caption itself has
% its own 0.3\baselineskip top skip and therefore does not need extra glue.
\def\examfigure@emitinline{%
    % Marker figures are already framed as a continuation slice.  Do not add
    % the ordinary inline-figure glue here, otherwise the slice acquires the
    % same visible seam as a float separated by \intextsep.
    \ifdefined\ifexam@marker@floatcontext
        \ifexam@marker@floatcontext\else\vspace{0.45\baselineskip}\fi
    \else
        \vspace{0.45\baselineskip}%
    \fi
    \ifdefined\ifexam@marker@floatcontext
        \ifexam@marker@floatcontext
            \begin{tcolorbox}[examanswermarkerfigure]%
                \box\examfigurebox
            \end{tcolorbox}%
        \else
            \box\examfigurebox
        \fi
    \else
        \box\examfigurebox
    \fi
    \par
    \ifdefined\ifexam@marker@floatcontext
        \ifexam@marker@floatcontext\else\vspace{0.4\baselineskip}\fi
    \else
        \vspace{0.4\baselineskip}%
    \fi
}
\def\examfigure@emitfloatwithspec#1#2{%
    \ifdim\examfigureheight>\textheight
        \PackageError{styles}{Figure box too tall for one page}{The figure #2 cannot fit without shrinking; fix the source figure.}%
        \examfigure@emitinline
    \else
        \global\@topnum\z@
        \examfigure@recordsourcetrue
        \begin{figure}[#1]
            \centering
            \ifdefined\ifexam@marker@floatcontext
                \ifexam@marker@floatcontext
                    \begin{tcolorbox}[examanswermarkerfigure]%
                        \box\examfigurebox
                    \end{tcolorbox}%
                \else
                    \box\examfigurebox
                \fi
            \else
                \box\examfigurebox
            \fi
        \end{figure}%
    \fi
}
\def\examfigure@emitfloat#1{\examfigure@emitfloatwithspec{!ht}{#1}}
\def\examfigure@emittop#1{\examfigure@emitfloatwithspec{!t}{#1}}
\def\examfigure@queuepending{%
    \expandafter\ifx\csname examfigure@queued@\number\@currbox\endcsname\relax
        \expandafter\xdef\csname examfigure@queued@\number\@currbox\endcsname{1}%
        \expandafter\xdef\csname examfigure@reserved@\number\@currbox\endcsname
            {\the\dimexpr\ht\@currbox+\dp\@currbox+\textfloatsep\relax}%
        \global\advance\examfigure@pendingheight by \dimexpr\ht\@currbox+\dp\@currbox+\textfloatsep\relax
        % The output routine already reduces \@colroom when this top float is
        % inserted.  Do not pre-reserve the pending height in tcolorbox: that
        % double-counts the source-page float and strands the next answer box.
    \fi
}
\def\examfigure@queueconsumed{%
    \expandafter\ifx\csname examfigure@queued@\number\@currbox\endcsname\relax\else
        \expandafter\ifx\csname examfigure@consumed@\number\@currbox\endcsname\relax
            \expandafter\xdef\csname examfigure@consumed@\number\@currbox\endcsname{1}%
            \expandafter\ifx\csname examfigure@reserved@\number\@currbox\endcsname\relax\else
                \global\advance\examfigure@pendingheight by
                    -\csname examfigure@reserved@\number\@currbox\endcsname\relax
                \ifdim\examfigure@pendingheight<0pt
                    \global\examfigure@pendingheight=0pt\relax
                \fi
                \ifnum\examfigure@currenttoppage=\c@page\relax\else
                    \global\examfigure@currenttopheight=0pt\relax
                    \global\examfigure@currenttoppage=\c@page\relax
                \fi
                \global\advance\examfigure@currenttopheight by
                    \csname examfigure@reserved@\number\@currbox\endcsname\relax
            \fi
        \fi
    \fi
}
% Attach the source page to the newly allocated float box.  This hook is
% deliberately local to exam figures; unrelated package floats keep LaTeX's
% normal placement behaviour.
\let\examfigure@orig@xfloat\@xfloat
\def\@xfloat#1[#2]{%
    \examfigure@orig@xfloat{#1}[#2]%
    \global\expandafter\let\csname examfigure@sourcepage@\number\@currbox\endcsname\relax
    \global\expandafter\let\csname examfigure@markertopfloat@\number\@currbox\endcsname\relax
    \ifdefined\ifexam@marker@floatcontext
        \ifexam@marker@floatcontext
            \global\expandafter\let\csname examfigure@markertopfloat@\number\@currbox\endcsname\@empty
        \fi
    \fi
    \ifexamfigure@recordsource
        \global\expandafter\let\csname examfigure@queued@\number\@currbox\endcsname\relax
        \global\expandafter\let\csname examfigure@reserved@\number\@currbox\endcsname\relax
        \global\expandafter\let\csname examfigure@consumed@\number\@currbox\endcsname\relax
        \expandafter\xdef\csname examfigure@sourcepage@\number\@currbox\endcsname{\number\examfigure@sourcepage}%
        \global\examfigure@recordsourcefalse
    \fi
}
% The output routine reads \textfloatsep only when a deferred top float is
% finally shipped.  By then the temporary marker-glue settings may already
% have been restored.  Track the last top float and suppress only the separator
% between a marker figure and its following answer-box slice; ordinary figures
% retain the document's normal \textfloatsep.
\newif\ifexamfigure@lasttopmarker
\let\examfigure@orig@comflelt\@comflelt
\def\examfigure@tracked@comflelt#1{%
    \expandafter\ifx\csname examfigure@markertopfloat@\number#1\endcsname\relax
        \examfigure@lasttopmarkerfalse
    \else
        \examfigure@lasttopmarkertrue
    \fi
    \examfigure@orig@comflelt{#1}%
}
\patchcmd{\@cflt}
  {\let\@elt\@comflelt}
  {\examfigure@lasttopmarkerfalse\let\@elt\examfigure@tracked@comflelt}
  {}
  {\PackageError{styles}{Failed to track the final top marker float}
    {The installed LaTeX output routine is incompatible with the marker figure layout guard.}}
\patchcmd{\@cflt}
  {\vskip\textfloatsep}
  {\ifexamfigure@lasttopmarker\else\vskip\textfloatsep\fi}
  {}
  {\PackageError{styles}{Failed to suppress marker top-float separation}
    {The installed LaTeX output routine is incompatible with the marker figure layout guard.}}
% On the source page, bypass only the top-float branch.  The ordinary bottom
% fallback (and the deferlist when no bottom placement is allowed) remains
% intact.  On the following page, the saved source page is strictly older,
% so the original top/bottom algorithm is restored.
\let\examfigure@orig@addtotoporbot\@addtotoporbot
\def\@addtotoporbot{%
    \expandafter\ifx\csname examfigure@sourcepage@\number\@currbox\endcsname\relax
        \examfigure@orig@addtotoporbot
    \else
        \ifnum\csname examfigure@sourcepage@\number\@currbox\endcsname<\c@page
            \examfigure@orig@addtotoporbot
            \if@insert
                \examfigure@queueconsumed
            \else
                \examfigure@queuepending
            \fi
        \else
            \@addtobot
            \if@insert\else\examfigure@queuepending\fi
        \fi
    \fi
}
\let\examfigure@orig@addtocurcol\@addtocurcol
\def\@addtocurcol{%
    \examfigure@orig@addtocurcol
    \expandafter\ifx\csname examfigure@sourcepage@\number\@currbox\endcsname\relax\else
        \if@insert\else
            \examfigure@queuepending
        \fi
    \fi
}
% A consumed mapped top float belongs to the page recorded in
% examfigure@currenttoppage.  queueconsumed starts a fresh height accumulator
% whenever output processing advances to another page.
\def\examfigure@dispatch#1#2#3{%
    \par\penalty0
    \examfigure@parseopts{#1}%
    \edef\examfigure@frozenprobnum{\theprobnum}%
    \examfigure@sourcepage=\c@page\relax
    \bitmappreparecaption{#2}%
    \edef\examfigure@call{\noexpand#3[\unexpanded\expandafter{\examfigure@opts}]{#2}}%
    \examfigure@buildbox{\examfigure@call}%
    \ifinner
        \ifx\examfigure@float\examfigure@modetop
            \PackageWarning{styles}{float=top ignored in inner mode for #2}%
        \fi
        \examfigure@emitinline
    \else
        \iftallpage
            \examfigure@emitinline
        \else\ifdim\pagegoal=\maxdimen
            \examfigure@emitinline
        \else\ifx\examfigure@float\examfigure@modeinline
            \examfigure@emitinline
        \else\ifx\examfigure@float\examfigure@modetop
            \examfigure@emittop{#2}%
        \else
            \examfigure@emitfloat{#2}%
        \fi\fi\fi\fi\fi
    \endgroup
}
\makeatother

\makeatletter
\newcommand{\bitmapfigure}[2][]{%
    \begingroup
    \examfigure@dispatch{#1}{#2}{\bitmapinclude}%
}

\newif\ifexamchoicegraphic
\newcount\examchoicegraphiccount

% 独立选择题选项图，供 choices/choicesfive 承载；尺寸由调用方控制。
% 全局标记只供 \choices 测量阶段识别图片选项；每次测量前都会重置。
\newcommand{\choicebitmap}[2][]{%
    \global\examchoicegraphictrue
    \begingroup
    \adjustbox{#1,max width=\linewidth,valign=c}{\includegraphics{#2}}%
    \endgroup
}

\newcommand{\choicetexfigure}[2][]{%
    \global\examchoicegraphictrue
    \begingroup
    \adjustbox{#1,max width=\linewidth,valign=c}{%
        \begingroup
        \texboxparagraphreset
        \texboxtikzfontredirect
        \input{#2}%
        \endgroup
    }%
    \endgroup
}

\newcommand{\texfigure}[2][]{%
    \begingroup
    \examfigure@dispatch{#1}{#2}{\texinclude}%
}

% Controlled composition interface for one logical figure made from several
% bitmap/TeX panels. Low-level bitmapinclude/texinclude are allowed only in
% this body; the outer dispatch keeps ordinary examfigure placement/caption
% behaviour and avoids ad-hoc figure environments in正文.
\def\examfiguregroup@body{}
\newcommand{\examfiguregroup@render}[2][]{%
    \texboxparagraphreset
    \examfiguregroup@body
}
\NewDocumentCommand{\examfiguregroup}{O{} m +m}{%
    \begingroup
    \def\examfiguregroup@body{#3}%
    \examfigure@dispatch{#1}{#2}{\examfiguregroup@render}%
}

% Narrow controlled exception for figures whose source is a page crop from a
% PDF. Graphicx keys stay on the inner image; float/placement keys are
% consumed by examfigure.
\newcommand{\sourcefigure@render}[2][]{%
    \adjustbox{valign=c,max width=\linewidth}{\originalincludegraphics[#1]{#2}}%
}
\NewDocumentCommand{\sourcefigure}{O{} m}{%
    \begingroup
    \examfigure@dispatch{#1}{#2}{\sourcefigure@render}%
}
\makeatother

% 频率分布直方图纵轴标签：标签保持正文大小；密集柱高可用可选横移量错列，
% 引线仍准确落在对应柱高，避免靠缩小字号处理重叠。
\NewDocumentCommand{\frequencyylabel}{O{-4pt} m m m}{%
    \draw[black,line width=0.4pt]
        (axis cs:#2,#3) -- ([xshift=#1]axis cs:#2,#3);
    \node[anchor=east,inner sep=1pt,font=\normalsize]
        at ([xshift=#1]axis cs:#2,#3) {#4};
}

% 列表与枚举
\usepackage{enumitem}
\AtBeginEnvironment{center}{\problemcenterstart}
\AfterEndEnvironment{center}{\problemcenterend}
\AtBeginEnvironment{tabular}{\problemconsumeinlineblockprefix\examcapturetablefractionclearance\examtablefractioncontexttrue}
\AtBeginEnvironment{tabular*}{\examcapturetablefractionclearance\examtablefractioncontexttrue}

% 盒子宏包 (用于答案)
\usepackage[most]{tcolorbox}

% A top float reserves \@colroom even while \pagegoal still reflects the
% original page goal.  Breakable tcolorbox otherwise computes a box height
% larger than the space actually left in the column and can cover that float.
% Apply the reserve only on the first post-break calculation; ordinary page
% and floating/maxdimen paths retain their native height rules.
\makeatletter
\def\examfigure@tcb@normalheight{%
    \ifdim\pagegoal<\@colroom
        \tcbdimto\tcb@h@page{\pagegoal-\pagetotal+\tcb@compress@height}%
    \else
        \tcbdimto\tcb@h@page{\@colroom-\pagetotal+\tcb@compress@height}%
    \fi
}
\def\examfigure@tcb@reserveheight{%
    \tcbdimto\tcb@h@page{\textheight-\examfigure@pendingheight+\tcb@compress@height}%
    \global\examfigure@tcbreservefalse
    \global\examfigure@tcbreadyfalse
}
\def\examfigure@tcb@reserveheightfloat{%
    \dimen@=\dimexpr\textheight-\examfigure@pendingheight\relax
    \dimen@ii=\vsize
    \ifdim\pagetotal<0pt
        \advance\dimen@ii by \pagetotal
    \fi
    \ifdim\dimen@ii<0pt
        \dimen@ii=0pt
    \fi
    \ifdim\dimen@ii<\dimen@
        \tcbdimto\tcb@h@page{\dimen@ii}%
    \else
        \tcbdimto\tcb@h@page{\dimen@}%
    \fi
    \global\examfigure@tcbreservefalse
    \global\examfigure@tcbreadyfalse
}
\patchcmd{\tcb@comp@h@page}
  {\tcbdimto\tcb@h@page{\vsize}}
   {\ifexamfigure@tcbreserve
     \ifdim\examfigure@currenttopheight>0pt\relax
       \ifnum\examfigure@currenttoppage=\c@page\relax
         \examfigure@tcb@reserveheightfloat
       \else
         \tcbdimto\tcb@h@page{\vsize}
       \fi
     \else
     \ifdim\examfigure@pendingheight>0pt\relax
       \ifexamfigure@tcbready
         \examfigure@tcb@reserveheightfloat
       \else
         \ifnum\c@tcbbreakpart>0\relax
           \examfigure@tcb@reserveheightfloat
         \else
           \tcbdimto\tcb@h@page{\vsize}
         \fi
       \fi
     \else
       \tcbdimto\tcb@h@page{\vsize}
     \fi
     \fi
   \else
     \tcbdimto\tcb@h@page{\vsize}
   \fi}
  {}
  {\PackageError{styles}{Failed to patch tcolorbox floating page-height calculation}
    {The installed tcolorbox breakable implementation is incompatible with the figure layout guard.}}
\patchcmd{\tcb@comp@h@page}
  {\tcbdimto\tcb@h@page{\pagegoal-\pagetotal+\tcb@compress@height}}
  {\ifexamfigure@tcbreserve
     \ifdim\examfigure@currenttopheight>0pt\relax
       \ifnum\examfigure@currenttoppage=\c@page\relax
         \examfigure@tcb@reserveheightfloat
       \else
         \ifdim\examfigure@pendingheight>0pt\relax
           \examfigure@tcb@reserveheight
         \else
           \examfigure@tcb@normalheight
         \fi
       \fi
     \else
       \ifdim\examfigure@pendingheight>0pt\relax
         \ifexamfigure@tcbready
           \examfigure@tcb@reserveheight
         \else
           \ifnum\c@tcbbreakpart>0\relax
             \examfigure@tcb@reserveheight
           \else
             \examfigure@tcb@normalheight
           \fi
         \fi
       \else
         \examfigure@tcb@normalheight
       \fi
     \fi
   \else
     \examfigure@tcb@normalheight
   \fi}
  {}
  {\PackageError{styles}{Failed to patch tcolorbox page-height calculation}
  {The installed tcolorbox breakable implementation is incompatible with the figure layout guard.}}
\patchcmd{\tcb@split@pagebreak}
  {\iftcb@multicol}
  {\global\examfigure@tcbreadytrue\iftcb@multicol}
  {}
  {\PackageError{styles}{Failed to patch tcolorbox page-break reserve hook}
    {The installed tcolorbox breakable implementation is incompatible with the figure layout guard.}}
\makeatother

\usepackage{ulem}
\newcommand{\examblankrule}{\rule[-0.2ex]{5em}{0.4pt}}
\newcommand{\fillinblank}{%
    \ifmmode
        \mathord{\text{\examblankrule}}%
    \else
        \unskip\nobreak\examblankrule%
    \fi
}

% ==============================================================================
% 4. 核心命令：题目、选项
% ==============================================================================

% --- A. 题目计数器 ---
% 绑定到 chapter，每次新试卷(chapter)自动重置为 0
% 注意：虽然 chapter 设置了 numbering=false，但计数器依然在后台运作
\newcounter{probnum}[chapter]
\newlength{\problembodyindent}
\setlength{\problembodyindent}{1.5em}
\newif\ifproblembodyactive
\newif\ifproblemfirstparagraph
\newif\ifproblemheadpending
\newcommand{\problemnumbersuffix}{}
\newcommand{\reuseproblemnumber}{%
    \addtocounter{probnum}{-1}%
}

\newcommand{\problemresetlayout}{%
    \leftskip=0pt\relax
    \hangindent=0pt\relax
    \hangafter=1\relax
    \parshape=0\relax
}

\newcommand{\problemheadbox}{%
    \makebox[\problembodyindent][r]{\textbf{\theprobnum\problemnumbersuffix.}\ }%
}

\newcommand{\problemheadprefix}{%
    \textbf{\theprobnum\problemnumbersuffix.}\ %
}
\newcommand{\problemoriginallabel}{}%

\newcommand{\problembodyfinish}{%
    \ifhmode\par\fi
    \global\problembodyactivefalse
    \global\problemfirstparagraphfalse
    \global\problemheadpendingfalse
    \everypar{}%
    \parfillskip=0pt plus 1fil\relax
    \problemresetlayout
}
\let\problembodyend\problembodyfinish

\newcommand{\problemliststart}{%
    \examproseparagraph
    \ifproblembodyactive
        \ifhmode\unskip\par\fi
        \ifproblemheadpending
            \global\problemfirstparagraphfalse
            \let\problemoriginallabel\makelabel
            \def\makelabel##1{%
                \ifproblemheadpending
                    \global\problemheadpendingfalse
                    \llap{\problemheadbox\hspace{0.5em}}%
                \fi
                \problemoriginallabel{##1}%
            }%
        \fi
    \fi
}

\newcommand{\problemblockstart}{%
    \examproseparagraph
    \ifproblembodyactive
        \ifhmode\unskip\par\fi
        \ifproblemheadpending
            \global\problemheadpendingfalse
            \global\problemfirstparagraphfalse
            \llap{\problemheadprefix}%
        \fi
    \fi
}

\newcommand{\problemconsumeinlineblockprefix}{%
    \examproseparagraph
    \ifproblemheadpending
        \global\problemheadpendingfalse
        \leavevmode\llap{\problemheadprefix}%
    \fi
}

\newcommand{\problemcenterstart}{%
    \examproseparagraph
    \ifproblembodyactive
        \ifhmode\unskip\par\fi
        \ifproblemheadpending
            \global\problemfirstparagraphfalse
            \linewidth=\dimexpr\linewidth-\problembodyindent\relax
            \hsize=\linewidth
        \fi
    \fi
}

\newcommand{\problemcenterend}{%
    \ifproblembodyactive
        \global\everypar{\problemcontinuationpar}%
    \fi
}

\makeatletter
\patchcmd{\@endparenv}
  {\addpenalty\@endparpenalty\addvspace\@topsepadd\@endpetrue}
  {\addpenalty\@endparpenalty\addvspace\@topsepadd%
   \ifproblembodyactive
     \global\@endpefalse
     \global\everypar{\problemcontinuationpar}%
   \else
     \@endpetrue
   \fi}
  {}{}

\newcommand{\problemcontinuationpar}{%
    \examproseparagraph
    \ifproblembodyactive
        \ifnum\@listdepth=\z@
            \ifproblemfirstparagraph
                \leftskip=0pt\relax
                \hangindent=\problembodyindent
                \hangafter=1
                \parshape=0\relax
                \ifproblemheadpending
                    \problemheadbox
                    \global\problemheadpendingfalse
                \fi
                \global\problemfirstparagraphfalse
            \else
                \leftskip=\problembodyindent\relax
                \hangindent=0pt\relax
                \hangafter=1\relax
                \parshape=0\relax
            \fi
        \fi
    \fi
}
\makeatother

% --- B. 题目命令 \problem ---
\newcommand{\problem}[1][]{%
    \problembodyfinish
    \examavoidshortpage
    \stepcounter{probnum}%
    \def\problemnumbersuffix{#1}%
    \global\problembodyactivetrue
    \global\problemfirstparagraphtrue
    \global\problemheadpendingtrue
    \examproseparagraph
    \everypar{\problemcontinuationpar}%
    \ignorespaces%
}
\let\endproblem\problembodyfinish

% --- C. 智能选项 \choices ---
\newcommand{\choiceentry}[2]{%
    \hangindent=1.8em\hangafter=1\noindent\textbf{#1.}~#2%
}

\newcommand{\choicecell}[3]{%
    \parbox[t]{#1}{\raggedright\textbf{#2.}~#3}%
}

% 四个图片选项专用：复用普通四选项的 0.22\examchoicewidth 列和 hfill，
% 保证 A/B/C/D 的列起点一致；标签与图片都按中心基线排放。
\newcommand{\choicegraphiccell}[3]{%
    \makebox[#1][l]{%
        \adjustbox{valign=c}{\textbf{#2.}}\hspace{0.25em}%
        #3%
    }%
}

\newdimen\examchoicewidth
\newdimen\examchoicewidthA
\newdimen\examchoicewidthB
\newdimen\examchoicewidthC
\newdimen\examchoicewidthD
\newdimen\examchoicewidthE
\newdimen\examchoicemaxwidth

\newcommand{\choices}[4]{%
    \problembodyend
    % Keep the transition mildly preferred but breakable; do not lock the
    % entire question to the current page when only the stem's first line fits.
    \par\penalty300
    \vspace{4pt}
    \begingroup
    \setlength{\leftskip}{\problembodyindent}
    \examchoicewidth=\dimexpr\linewidth-\problembodyindent\relax
    \examchoicegraphiccount=0
    \global\examchoicegraphicfalse
    \settowidth{\examchoicewidthA}{\textbf{A.}~#1}
    \ifexamchoicegraphic\advance\examchoicegraphiccount by 1\relax\fi
    \global\examchoicegraphicfalse
    \settowidth{\examchoicewidthB}{\textbf{B.}~#2}
    \ifexamchoicegraphic\advance\examchoicegraphiccount by 1\relax\fi
    \global\examchoicegraphicfalse
    \settowidth{\examchoicewidthC}{\textbf{C.}~#3}
    \ifexamchoicegraphic\advance\examchoicegraphiccount by 1\relax\fi
    \global\examchoicegraphicfalse
    \settowidth{\examchoicewidthD}{\textbf{D.}~#4}
    \ifexamchoicegraphic\advance\examchoicegraphiccount by 1\relax\fi
    \examchoicemaxwidth=\examchoicewidthA
    \ifdim\examchoicewidthB>\examchoicemaxwidth \examchoicemaxwidth=\examchoicewidthB \fi
    \ifdim\examchoicewidthC>\examchoicemaxwidth \examchoicemaxwidth=\examchoicewidthC \fi
    \ifdim\examchoicewidthD>\examchoicemaxwidth \examchoicemaxwidth=\examchoicewidthD \fi
    \addtolength{\examchoicemaxwidth}{2em}

    \vspace{-0.2em}\noindent
    \ifnum\examchoicegraphiccount=4
        \choicegraphiccell{\dimexpr0.22\examchoicewidth\relax}{A}{#1}\hfill
        \choicegraphiccell{\dimexpr0.22\examchoicewidth\relax}{B}{#2}\hfill
        \choicegraphiccell{\dimexpr0.22\examchoicewidth\relax}{C}{#3}\hfill
        \choicegraphiccell{\dimexpr0.22\examchoicewidth\relax}{D}{#4}%
    \else\ifdim\examchoicemaxwidth<0.22\examchoicewidth
        \choicecell{\dimexpr0.22\examchoicewidth\relax}{A}{#1}\hfill
        \choicecell{\dimexpr0.22\examchoicewidth\relax}{B}{#2}\hfill
        \choicecell{\dimexpr0.22\examchoicewidth\relax}{C}{#3}\hfill
        \choicecell{\dimexpr0.22\examchoicewidth\relax}{D}{#4}%
    \else
        \ifdim\examchoicemaxwidth<0.47\examchoicewidth
            \choicecell{\dimexpr0.47\examchoicewidth\relax}{A}{#1}\hfill
            \choicecell{\dimexpr0.47\examchoicewidth\relax}{B}{#2}\par
            \vspace{0.2em}
            \choicecell{\dimexpr0.47\examchoicewidth\relax}{C}{#3}\hfill
            \choicecell{\dimexpr0.47\examchoicewidth\relax}{D}{#4}%
        \else
            \choiceentry{A}{#1}\par
            \choiceentry{B}{#2}\par
            \choiceentry{C}{#3}\par
            \choiceentry{D}{#4}%
        \fi
    \fi\fi
    \global\examchoicegraphicfalse
    \par
    \endgroup
}

% --- D. 五选项 choicesfive ---
% 排布：5*1行 > 3+2 > 2+2+1 > 竖排
% 对齐规则：D/E 对齐第一行的 A/B（3+2 时），E 对齐 A（2+2+1 时）
\newcommand{\choicesfive}[5]{%
    \problembodyend
    \par\penalty300
    \vspace{4pt}
    \begingroup
    \setlength{\leftskip}{\problembodyindent}
    \examchoicewidth=\dimexpr\linewidth-\problembodyindent\relax
    \settowidth{\examchoicewidthA}{\textbf{A.}~#1}
    \settowidth{\examchoicewidthB}{\textbf{B.}~#2}
    \settowidth{\examchoicewidthC}{\textbf{C.}~#3}
    \settowidth{\examchoicewidthD}{\textbf{D.}~#4}
    \settowidth{\examchoicewidthE}{\textbf{E.}~#5}
    \examchoicemaxwidth=\examchoicewidthA
    \ifdim\examchoicewidthB>\examchoicemaxwidth \examchoicemaxwidth=\examchoicewidthB \fi
    \ifdim\examchoicewidthC>\examchoicemaxwidth \examchoicemaxwidth=\examchoicewidthC \fi
    \ifdim\examchoicewidthD>\examchoicemaxwidth \examchoicemaxwidth=\examchoicewidthD \fi
    \ifdim\examchoicewidthE>\examchoicemaxwidth \examchoicemaxwidth=\examchoicewidthE \fi
    \addtolength{\examchoicemaxwidth}{2em}
    \vspace{-0.2em}\noindent
    % 5 per row
    \ifdim\examchoicemaxwidth<0.18\examchoicewidth
        \choicecell{\dimexpr0.18\examchoicewidth\relax}{A}{#1}\hfill
        \choicecell{\dimexpr0.18\examchoicewidth\relax}{B}{#2}\hfill
        \choicecell{\dimexpr0.18\examchoicewidth\relax}{C}{#3}\hfill
        \choicecell{\dimexpr0.18\examchoicewidth\relax}{D}{#4}\hfill
        \choicecell{\dimexpr0.18\examchoicewidth\relax}{E}{#5}%
    \else
    % 3+2: row1=A/B/C, row2=D/E (D对齐A, E对齐B)
    \ifdim\examchoicemaxwidth<0.30\examchoicewidth
        \choicecell{\dimexpr0.30\examchoicewidth\relax}{A}{#1}\hfill
        \choicecell{\dimexpr0.30\examchoicewidth\relax}{B}{#2}\hfill
        \choicecell{\dimexpr0.30\examchoicewidth\relax}{C}{#3}\par
        \vspace{0.2em}
        \choicecell{\dimexpr0.30\examchoicewidth\relax}{D}{#4}\hfill
        \choicecell{\dimexpr0.30\examchoicewidth\relax}{E}{#5}\hfill
        \makebox[\dimexpr0.30\examchoicewidth\relax][l]{}%
    \else
    % 2+2+1: row1=A/B, row2=C/D, row3=E (E对齐A)
    \ifdim\examchoicemaxwidth<0.47\examchoicewidth
        \choicecell{\dimexpr0.47\examchoicewidth\relax}{A}{#1}\hfill
        \choicecell{\dimexpr0.47\examchoicewidth\relax}{B}{#2}\par
        \vspace{0.2em}
        \choicecell{\dimexpr0.47\examchoicewidth\relax}{C}{#3}\hfill
        \choicecell{\dimexpr0.47\examchoicewidth\relax}{D}{#4}\par
        \vspace{0.2em}
        \choiceentry{E}{#5}%
    \else
    % vertical
        \choiceentry{A}{#1}\par
        \choiceentry{B}{#2}\par
        \choiceentry{C}{#3}\par
        \choiceentry{D}{#4}\par
        \choiceentry{E}{#5}%
    \fi\fi\fi
    \par
    \endgroup
}

% ==============================================================================
% 5. 答案控制系统 (框式设计)
% ==============================================================================
\newif\ifshowanswer
\showanswertrue     % 【默认开启】

% 答案框内的 enumerate 不是悬挂列表，而是带自动标签的普通段落流。
% 第一项可紧跟框标题；嵌套第一项可紧跟父项标签。其余项目及项目内
% 的显式新段落均沿用答案正文的 2em 首行缩进，续行回到框体左边界。
\makeatletter
\newif\ifexamanswerenumfirstitem
\newif\ifexamanswerenumcontinuefirst

\newcommand{\examanswerenumemititem}[1]{%
    \ifexamanswerenumfirstitem
        \ifexamanswerenumcontinuefirst
            \unskip
        \else
            \par\indent
        \fi
    \else
        \par\indent
    \fi
    \examanswerenumfirstitemfalse
    #1\ignorespaces
}

\newcommand{\examanswerenumstep}{%
    % Hyperref inserts a zero-width destination between adjacent generated
    % labels; xeCJK then treats that node as a punctuation boundary and adds
    % visible space in `（2）（i）`.  The saved kernel refstepcounter retains
    % counter/\label semantics without injecting that horizontal node.
    \@ifundefined{H@refstepcounter}{%
        \refstepcounter{\@enumctr}%
    }{%
        \H@refstepcounter{\@enumctr}%
    }%
}

\newcommand{\examanswerenumitem}{%
    \@ifnextchar[\examanswerenumitemexplicit\examanswerenumitemautomatic
}

\newcommand{\examanswerenumitemautomatic}{%
    \examanswerenumstep
    \expandafter\examanswerenumemititem\expandafter{%
        \csname label\@enumctr\endcsname
    }%
}

\def\examanswerenumitemexplicit[#1]{%
    \examanswerenumstep
    \examanswerenumemititem{#1}%
}

\newenvironment{examanswerenumerate}[1][]{%
    \ifnum\@enumdepth>3\relax
        \@toodeep
    \else
        \ifhmode
            \examanswerenumcontinuefirsttrue
        \else
            \examanswerenumcontinuefirstfalse
        \fi
        \advance\@enumdepth\@ne
        \edef\@enumctr{enum\romannumeral\the\@enumdepth}%
        \setcounter{\@enumctr}{0}%
        \examanswerenumfirstitemtrue
        \let\item\examanswerenumitem
    \fi
    \ignorespaces
}{%
    \ifnum\@enumdepth=1\relax
        \par
    \fi
}

\newcommand{\examanswerenumeratecontext}{%
    \def\labelenumi{（\arabic{enumi}）}%
    \def\labelenumii{（\roman{enumii}）}%
    % Paragraph-style items temporarily replace \item.  LaTeX's center,
    % flush and quote environments are based on trivlist and still need the
    % kernel \item internally, so restore it only inside those local groups.
    \let\examanswerstandarditem\item
    \let\examanswerstandardtrivlist\trivlist
    \def\trivlist{\let\item\examanswerstandarditem\examanswerstandardtrivlist}%
    \let\enumerate\examanswerenumerate
    \let\endenumerate\endexamanswerenumerate
}
\makeatother

% 定义答案盒子的样式
\tcbset{
    examanswer/.style={
        enhanced jigsaw,
        breakable,             % 允许跨栏
        colback=gray!5,        % 极淡的灰色背景
        colframe=gray!40,      % 边框颜色
        boxrule=0.5pt,         % 细边框
        arc=1mm,               % 轻微圆角
        left=3pt,right=2pt, top=1pt, bottom=1pt,
        before skip=0.3em,
        after skip=0.3em,
        before upper={\examwidoworphanrules\examproseparagraph\setlength{\parindent}{2em}\examanswerenumeratecontext\tolerance=3200\emergencystretch=8em\setlength{\jot}{4pt}\lineskiplimit=3pt\lineskip=4pt\selectfont},
        fontupper=\zihao{-4},
    },
    % Marker slices are named styles rather than inline key lists: keyval
    % normalizes spaces in dynamically supplied lists, while these definitions
    % retain tcolorbox's `before skip`/`after skip` keys verbatim.
    examanswermarkerfirst/.style={
        examanswer,
        bottomrule=0pt,
        arc=0pt,
        after skip=0pt,
    },
    examanswermarkermiddle/.style={
        examanswer,
        toprule=0pt,
        bottomrule=0pt,
        sharp corners,
        before skip=0pt,
        after skip=0pt,
    },
    examanswermarkerlast/.style={
        examanswer,
        toprule=0pt,
        arc=0pt,
        before skip=0pt,
        boxsep=0.5pt,
        top=0.5pt,
        bottom=0.5pt,
        after skip=0.1em,
    },
    examanswermarkerfigure/.style={
        examanswer,
        unbreakable,
        toprule=0pt,
        bottomrule=0pt,
        sharp corners,
        before skip=0pt,
        after skip=0pt,
    },
}

% 答案/解析环境；旧命令写法 \answer{...}、\solution{...} 继续兼容。
\newbox\examhiddenanswerbox
\makeatletter
\newcommand{\examallownextpagebreak}{\@nobreakfalse}
\newcommand{\answerbegin}{%
    \problembodyend
    \ifshowanswer
        \examallownextpagebreak
        \begin{tcolorbox}[examanswer]
            \textbf{【答案】}\ignorespaces
    \else
        \setbox\examhiddenanswerbox=\vbox\bgroup\hsize=\linewidth
    \fi
}
\newcommand{\answerend}{%
    \ifshowanswer
        \end{tcolorbox}
    \else
        \egroup
    \fi
}
\NewDocumentCommand{\answer}{+g}{%
    \IfNoValueTF{#1}{\answerbegin}{\answerbegin #1\answerend}%
}
\let\endanswer\answerend

\newcommand{\solutionbegin}{%
    \problembodyend
    \ifshowanswer
        \examallownextpagebreak
        \begin{tcolorbox}[examanswer]
            \textbf{【解析】}\ignorespaces
    \else
        \setbox\examhiddenanswerbox=\vbox\bgroup\hsize=\linewidth
    \fi
}
\newcommand{\solutionend}{%
    \ifshowanswer
        \end{tcolorbox}
    \else
        \egroup
    \fi
}
\makeatother

% Explicit solution/answer figure marker.  It is consumed only by the body
% collector below; outside a collected body its error keeps accidental use
% visible instead of silently changing layout.
\NewDocumentCommand{\solutionfigure}{+m}{%
    \PackageError{styles}{solutionfigure used outside answer/solution body}%
        {Wrap the complete figure command in an answer or solution body.}%
}

\makeatletter
\ExplSyntaxOn
\def\exam@marker@stop{\relax}
\newif\ifexam@marker@first
\newif\ifexam@marker@forcefirst
\newif\ifexam@marker@hasmarker
\newif\ifexam@marker@floatcontext
\exam@marker@floatcontextfalse
\newif\ifexam@marker@floatgluezero
\exam@marker@floatgluezerofalse
\newskip\exam@marker@savedintextsep
\newskip\exam@marker@savedtextfloatsep
\newskip\exam@marker@savedfloatsep
% A marker float must remain a real float, but its external float glue would
% visibly break the surrounding examanswer frame.  Keep the original lengths
% until FloatBarrier has placed all marker floats, then restore them before
% subsequent ordinary figures are emitted.
\def\exam@marker@beginfloatglue{%
    \ifexam@marker@floatgluezero\else
        \global\exam@marker@savedintextsep=\intextsep
        \global\exam@marker@savedtextfloatsep=\textfloatsep
        \global\exam@marker@savedfloatsep=\floatsep
        \global\intextsep=0pt
        \global\textfloatsep=0pt
        \global\floatsep=0pt
        \global\exam@marker@floatgluezerotrue
    \fi
}
\def\exam@marker@endfloatglue{%
    \ifexam@marker@floatgluezero
        \global\intextsep=\exam@marker@savedintextsep
        \global\textfloatsep=\exam@marker@savedtextfloatsep
        \global\floatsep=\exam@marker@savedfloatsep
        \global\exam@marker@floatgluezerofalse
    \fi
}
% A blank line between two solutionfigure markers is tokenized as \par while
% the body is collected.  Strip only paragraph tokens and surrounding spaces
% for the emptiness test; any real text, math, or layout command remains a
% non-blank segment and is rendered normally.
\prg_new_conditional:Npnn \exam_marker_if_segment_blank:n {p,T,F,TF}
  {
    \tl_set:Nn \l_tmpa_tl {##1}
    \tl_remove_all:Nn \l_tmpa_tl {\par}
    \tl_if_blank:VTF \l_tmpa_tl
      {\prg_return_true:}
      {\prg_return_false:}
  }
% A solution figure is a real float and therefore cannot live inside one
% physical breakable tcolorbox.  The collector below emits visual *slices*
% of one logical box instead: the slice before a float has no bottom rule or
% lower rounding, the slice after it has no top rule or upper rounding, and
% any intervening slice has neither.  The side rules and shared background
% continue through the float's source page, so this preserves the float's
% placement/back-fill behaviour without exposing separate rounded boxes.
\long\def\exam@marker@rendersegment#1#2#3{%
    \ifexam@marker@forcefirst
        \begin{tcolorbox}[#3]%
            \textbf{#1}\ignorespaces
            #2
        \end{tcolorbox}%
        \global\exam@marker@firstfalse
        \global\exam@marker@forcefirstfalse
    \else
      \exam_marker_if_segment_blank:nTF {#2}{%
        % Empty source segments (including one or more \par tokens) do not
        % create a visual tcolorbox slice.
      }{%
        \begin{tcolorbox}[#3]%
            \ifexam@marker@first
                \textbf{#1}\ignorespaces
            \fi
            #2
        \end{tcolorbox}%
        \global\exam@marker@firstfalse
      }%
    \fi
}
% The figure command itself still owns the standard float placement.  Mark
% its construction, so the float box can receive the same gray background and
% side rules as the surrounding text slices without nesting a float inside a
% tcolorbox (which would disable floating in inner mode).
\def\exam@marker@emitfigure#1{%
    \exam@marker@beginfloatglue
    \global\exam@marker@floatcontexttrue
    #1%
    \global\exam@marker@floatcontextfalse
}
\long\def\exam@marker@parse#1\solutionfigure#2#3\exam@marker@end{%
    \ifx#2\exam@marker@stop
        \ifexam@marker@hasmarker
            % The final slice follows the last float.  Its top edge is a
            % continuation edge; only its bottom edge closes the container.
            \exam@marker@rendersegment{\exam@marker@heading}{#1}{examanswermarkerlast}%
        \else
            \exam@marker@rendersegment{\exam@marker@heading}{#1}{examanswer}%
        \fi
    \else
        \global\exam@marker@hasmarkertrue
        \ifexam@marker@first
            \global\exam@marker@forcefirsttrue
            % First slice: keep the top rounding, but leave the lower edge
            % open for the continuation slice after the float.
            \exam@marker@rendersegment{\exam@marker@heading}{#1}{examanswermarkerfirst}%
        \else
            % Middle slice between two floats: no horizontal seam at all.
            \exam@marker@rendersegment{\exam@marker@heading}{#1}{examanswermarkermiddle}%
        \fi
        \exam@marker@emitfigure{#2}%
        % Pass the explicit end sentinel into each recursive parse.  Without
        % it, the final segment would scan forward to a later solutionfigure
        % in the document and be rendered a second time.
        \exam@marker@parse#3\solutionfigure{\exam@marker@stop}\exam@marker@end
    \fi
}
\long\def\exam@marker@rendersliced#1#2{%
    \def\exam@marker@heading{#1}%
    \exam@marker@firsttrue
    \exam@marker@forcefirstfalse
    \exam@marker@hasmarkerfalse
    \exam@marker@parse#2\solutionfigure{\exam@marker@stop}\exam@marker@end
    \ifexam@marker@hasmarker
        \FloatBarrier
        \exam@marker@endfloatglue
    \fi
}
% A source enumerate may span one or more solutionfigure markers.  Splitting
% such a body into physical tcolorbox slices would close a slice while the
% environment group is still open.  Keep that body in one breakable box and
% let its marked figures use bitmapfigure/texfigure's existing inner-mode
% block rendering.  Marker bodies without enumerate retain real-float slices.
\long\def\exam@marker@renderenumerated#1#2{%
    \begin{tcolorbox}[examanswer]%
        \textbf{#1}\ignorespaces
        \long\def\solutionfigure##1{%
            \begingroup
                \parindent=0pt\relax
                ##1%
            \endgroup
        }%
        #2
    \end{tcolorbox}%
}
\long\def\exam@marker@render#1#2{%
    \ifshowanswer
        \problembodyend
        \examallownextpagebreak
        \tl_if_in:nnTF {#2} {\begin{enumerate}}%
            {\tl_if_in:nnTF {#2} {\solutionfigure}%
                {\exam@marker@renderenumerated{#1}{#2}}%
                {\exam@marker@rendersliced{#1}{#2}}}%
            {\exam@marker@rendersliced{#1}{#2}}%
    \fi
}
\long\def\exam@collectsolution#1{\exam@marker@render{\textbf{【解析】}}{#1}}
\long\def\exam@collectanswer#1{\exam@marker@render{\textbf{【答案】}}{#1}}
\NewDocumentCommand{\solution}{+g}{%
    \IfNoValueTF{#1}{\Collect@Body\exam@collectsolution}{\exam@marker@render{\textbf{【解析】}}{#1}}%
}
\RenewDocumentCommand{\answer}{+g}{%
    \IfNoValueTF{#1}{\Collect@Body\exam@collectanswer}{\exam@marker@render{\textbf{【答案】}}{#1}}%
}
\let\endanswer\relax
\ExplSyntaxOff
\makeatother

% ==============================================================================
% 5.1 解析区行间公式归一化宏
% ==============================================================================
\NewDocumentCommand{\examdisplayaligned}{o +m}{%
    \[
    \IfValueTF{#1}{\begin{aligned}[#1]}{\begin{aligned}}
    #2
    \end{aligned}
    \]
}
\NewDocumentCommand{\examdisplaycases}{+m +m}{%
    \[
    #1\begin{cases}
    #2
    \end{cases}
    \]
}
\NewDocumentCommand{\examdisplayarray}{+m m +m +m}{%
    \[
    \begingroup
    \examcapturetablefractionclearance
    \examtablefractioncontexttrue
    #1\begin{array}{#2}
    #3
    \end{array}#4
    \endgroup
    \]
}
\NewDocumentCommand{\examdisplaychain}{+m}{%
    \begin{align*}
    #1
    \end{align*}
}

% ==============================================================================
% 6. 图表标题定制 (第X题图)
% ==============================================================================
\DeclareCaptionLabelFormat{examfig}{第\theprobnum 题图}
\DeclareCaptionLabelFormat{examtab}{第\theprobnum 题表}

\captionsetup[figure]{
    labelformat=examfig,
    labelsep=none,
    font=small,
    aboveskip=.3\baselineskip,
    belowskip=0pt
}
\captionsetup[table]{
    labelformat=examtab,
    labelsep=none,
    font=small
}

% ==============================================================================
% 7. 辅助符号与工具
% ==============================================================================
\renewcommand{\thefootnote}{\circled{\arabic{footnote}}}
\newcommand{\circled}[2][]{\tikz[baseline=(char.base)]
    {\node[shape=circle,draw,inner sep=0.8pt](char){\small #2};}}
\makeatletter
\renewcommand{\tagform@}[1]{\maketag@@@{\circled{#1}}}
\makeatother

\newcommand{\R}{\symbfup{R}}
\newcommand{\Z}{\symbfup{Z}}
\newcommand{\N}{\mathbb{N}}
\renewcommand{\Pr}{P}
\renewcommand{\d}{\text{d}}
\newcommand{\e}{\text{e}}
\newcommand{\degree}{^\circ}
\newcommand{\bs}[1]{\symbfit{#1}}
\let\oldfrac\frac
\newcommand{\examtablefractionpad}[1]{%
    \ifexamtablefractioncontext
        \mathinner{\raisebox{0pt}%
            [\dimexpr\height+\examtablefractionclearance\relax]%
            [\dimexpr\depth+\examtablefractionclearance\relax]%
            {\ensuremath{#1}}}%
    \else
        #1%
    \fi
}
\renewcommand{\frac}[2]{%
    \mathchoice
        {\examtablefractionpad{\dfrac{#1}{#2}}}      % 1. Display style；表格中保留大分式并补上下净空
        {\examtablefractionpad{\dfrac{#1}{#2}}}      % 2. Text style；行内仍强制使用大分式
        {\oldfrac{\scriptstyle #1}{\scriptstyle #2}} % 3. Script style (指数、上下标) -> 强制分子分母维持 scriptstyle
        {\oldfrac{\scriptscriptstyle #1}{\scriptscriptstyle #2}} % 4. Scriptscript style (更深层的下标) -> 恢复默认小分式
}

% 大运算符在行内公式中默认使用显示模式的格式（类似 \frac 的处理）
\makeatletter
% 辅助宏：sum 类大运算符（\mathchoice 控制大小，\limits 控制极限位置）
\newcommand{\make@display@bigop}[2]{% #1=备份命令, #2=原命令
  \renewcommand{#2}{\mathop{%
    \mathchoice
        {#1}%                              % 1. Display style -> 保持原样（已经是大符号）
        {\displaystyle #1}%                % 2. Text style (行内公式) -> 强制大符号
        {#1}%                              % 3. Script style (指数、上下标) -> 保持原样
        {#1}%                              % 4. Scriptscript style -> 保持原样
  }\limits}}
% 积分类必须把上下限直接挂在原生积分符上，保留字体的 MathKern 光学校正。
\newcommand{\apply@integral@scripts}[3]{%
  #1\IfValueT{#2}{_{#2}}\IfValueT{#3}{^{#3}}%
}
\newcommand{\make@display@bigint}[2]{% #1=备份命令, #2=原命令
  \RenewDocumentCommand{#2}{e{_^}}{%
    \mathchoice
        {\apply@integral@scripts{#1}{##1}{##2}}%              % Display style
        {\displaystyle\apply@integral@scripts{#1}{##1}{##2}}% Text style -> 强制大符号
        {\apply@integral@scripts{#1}{##1}{##2}}%              % Script style
        {\apply@integral@scripts{#1}{##1}{##2}}%              % Scriptscript style
  }}
\makeatother

% \let\oldlim\lim
% \renewcommand{\lim}{\oldlim\limits}

\newcommand{\myarc}[1]{\overset{\frown}{#1}}

\SetEnumitemKey{examhang}{%
    nosep,
    align=left,
    before=\problemliststart,
    labelindent=\problembodyindent,
    leftmargin=\dimexpr\problembodyindent+\labelwidth+\labelsep\relax,
    itemindent=0pt,
    listparindent=0pt,
    parsep=0pt,
    topsep=0pt,
    partopsep=0pt
}

\setlist[enumerate,1]{ % 仅针对第一层级
    examhang,
    label=（\arabic*）,
    labelsep=0em,
    widest*=1
}
\setlist[enumerate,2]{
    examhang,
    label=（\roman*）,
    labelsep=0em,
    labelindent=0pt,
    leftmargin=\dimexpr\labelwidth+\labelsep\relax,
    widest*=1
}

% ==============================================================================
% circlelist：智能多选题选项环境
%   对齐：标签（圈号）在左边距，正文缩进 \problembodyindent
%   布局：自动测量每条宽度，决定水平排列 / 双列 / 竖排
% ==============================================================================
\ExplSyntaxOn
\seq_new:N \g__circle_items_seq
\int_new:N \g__circle_count_int
\dim_new:N \g__circle_maxwidth_dim
\dim_new:N \g__circle_availwidth_dim
\box_new:N \g__circle_measure_box

\NewDocumentEnvironment{circlelist}{O{} +b}{
  \par\nopagebreak
  % 按 \item 拆分 body
  \regex_split:nnN { \c{item}\s* } { #2 } \l_tmpa_seq
  \seq_pop_left:NN \l_tmpa_seq \l_tmpa_tl
  \seq_gclear:N \g__circle_items_seq
  \seq_map_inline:Nn \l_tmpa_seq {
    \tl_set:Nn \l_tmpa_tl { ##1 }
    \tl_trim_spaces:N \l_tmpa_tl
    \seq_gput_right:NV \g__circle_items_seq \l_tmpa_tl
  }
  \int_gset:Nn \g__circle_count_int { \seq_count:N \g__circle_items_seq }
  \int_compare:nNnF { \g__circle_count_int } = { 0 } {
    % 测量每条宽度
    \dim_gzero:N \g__circle_maxwidth_dim
    \int_step_inline:nn { \g__circle_count_int } {
      \hbox_set:Nn \g__circle_measure_box {
        \circled{##1}~\seq_item:Nn \g__circle_items_seq { ##1 }
      }
      \dim_gset:Nn \g__circle_maxwidth_dim {
        \dim_max:nn { \g__circle_maxwidth_dim } { \box_wd:N \g__circle_measure_box }
      }
    }
    \dim_gadd:Nn \g__circle_maxwidth_dim { 1em }
    \dim_set:Nn \g__circle_availwidth_dim { \linewidth - \problembodyindent }
    \group_begin:
    \examproseparagraph
    \setlength{\leftskip}{\problembodyindent}
    \setlength{\parindent}{0pt}
    \dim_compare:nNnTF
      { \g__circle_maxwidth_dim * \g__circle_count_int }
      < { \g__circle_availwidth_dim }
    {
      % 全部在一行
      \dim_set:Nn \l_tmpa_dim { \g__circle_availwidth_dim / \g__circle_count_int }
      \int_step_inline:nn { \g__circle_count_int } {
        \makebox[\l_tmpa_dim][l]{
          \circled{##1}~\seq_item:Nn \g__circle_items_seq { ##1 }
        }
      }
    }{
      \dim_compare:nNnTF
        { \g__circle_maxwidth_dim * 2 }
        < { \g__circle_availwidth_dim }
      {
        % 每行两个
        \dim_set:Nn \l_tmpa_dim { ( \g__circle_availwidth_dim - 1em ) / 2 }
        \int_zero:N \l_tmpa_int
        \int_step_inline:nn { \g__circle_count_int } {
          \int_incr:N \l_tmpa_int
          \makebox[\l_tmpa_dim][l]{
            \circled{##1}~\seq_item:Nn \g__circle_items_seq { ##1 }
          }
          \int_compare:nNnTF { \l_tmpa_int } = { 2 } {
            \par\vspace{0.2em}
            \int_zero:N \l_tmpa_int
          }{
            \hfill
          }
        }
        \int_compare:nNnF { \l_tmpa_int } = { 0 } { \par }
      }{
        % 竖排
        \int_step_inline:nn { \g__circle_count_int } {
          \hangindent=1.8em\hangafter=1\noindent
          \circled{##1}~\seq_item:Nn \g__circle_items_seq { ##1 }
          \par
        }
      }
    }
    \par
    \group_end:
  }
  \par
}{}
\ExplSyntaxOff

% 修改公式与前后文的紧凑度
\makeatletter
\g@addto@macro\normalsize{
    % --- 控制行间公式 (Display Math) ---
	\setlength\abovedisplayskip{0.25em plus 0.5em minus 0.25em}
	\setlength\belowdisplayskip{0.25em plus 0.5em minus 0.25em}
	\setlength\abovedisplayshortskip{0em plus 2pt minus 0pt}
	\setlength\belowdisplayshortskip{0em plus 2pt minus 0pt}

    % --- 控制行内公式 (Inline Math) 的避让机制 ---
    % 设置阈值：当两行间距小于 1pt 时触发避让
    \setlength\lineskiplimit{3pt}
    % 设置避让间距：触发后强制留出 3pt 的缝隙，防止重叠
    \setlength\lineskip{3pt}
}
\makeatother

\makeatletter
\AtBeginDocument{
    % 处理 pi 为正体
    \let\oldpi\pi
    \renewcommand{\pi}{\symup\oldpi}

    % 组合数统一改成高中记号 C_n^m；无论源码使用 \binom / \dbinom / \tbinom，
    % 成稿都输出同一套样式。
    \RenewDocumentCommand{\binom}{m m}{\mathrm C_{#1}^{#2}}
    \RenewDocumentCommand{\dbinom}{m m}{\mathrm C_{#1}^{#2}}
    \RenewDocumentCommand{\tbinom}{m m}{\mathrm C_{#1}^{#2}}

    % 处理 nabla：不带 ^2 时加粗，带 ^2 时（Laplacian）保持原样
    \let\oldnabla\nabla
    \RenewDocumentCommand{\nabla}{e{^}}{%
        \IfValueTF{#1}{% 有上标
            \ifstrequal{#1}{2}{% 上标是 2
                \oldnabla^{#1}%
            }{% 上标不是 2
                \symbf{\oldnabla}^{#1}%
            }%
        }{% 无上标
            \symbf{\oldnabla}%
        }%
    }

    % --- 大运算符：在 \AtBeginDocument 中重定义，确保在 unicode-math 之后 ---
    % sum 类
    \let\oldsum\sum             \make@display@bigop{\oldsum}{\sum}
    \let\oldprod\prod           \make@display@bigop{\oldprod}{\prod}
    \let\oldcoprod\coprod       \make@display@bigop{\oldcoprod}{\coprod}
    \let\oldbigcup\bigcup       \make@display@bigop{\oldbigcup}{\bigcup}
    \let\oldbigcap\bigcap       \make@display@bigop{\oldbigcap}{\bigcap}
    \let\oldbigsqcup\bigsqcup   \make@display@bigop{\oldbigsqcup}{\bigsqcup}
    \let\oldbigvee\bigvee       \make@display@bigop{\oldbigvee}{\bigvee}
    \let\oldbigwedge\bigwedge   \make@display@bigop{\oldbigwedge}{\bigwedge}
    \let\oldbigoplus\bigoplus   \make@display@bigop{\oldbigoplus}{\bigoplus}
    \let\oldbigotimes\bigotimes \make@display@bigop{\oldbigotimes}{\bigotimes}
    \let\oldbigodot\bigodot     \make@display@bigop{\oldbigodot}{\bigodot}
    % 积分类
    \let\oldint\int             \make@display@bigint{\oldint}{\int}
    \let\oldiint\iint           \make@display@bigint{\oldiint}{\iint}
    \let\oldiiint\iiint         \make@display@bigint{\oldiiint}{\iiint}
    \let\oldiiiint\iiiint       \make@display@bigint{\oldiiiint}{\iiiint}
    \let\oldoint\oint           \make@display@bigint{\oldoint}{\oint}
}
\makeatother

% PDF 导航：目录与大纲均采用“年份 → 试卷”两级结构。
\usepackage[hidelinks,bookmarksdepth=chapter,hypertexnames=false]{hyperref}
\usepackage{bookmark}
\renewcommand{\i}{\ensuremath{\mathrm{i}}}
\newcommand{\examyear}[1]{%
    \clearpage
    \phantomsection
    \addcontentsline{toc}{part}{#1}%
}
